% ===== uphill/p2_c2.tex  (Problem 2, track 2) =====
% Written directly as a body file (not produced by tools/convert_cell.py).
% Title: Problem 2, Cell 2: $U(Q_5)$
% Body only: packages, theorem environments and macros are in preamble.tex.
% Labels are prefixed with "p2c2:".  Uses the framework of cell 1 (labels p2c1:*).
% Code: code/uphill_lower.py (claim L5), code/uphill_construct.py, code/uphill_eval.py.

\paragraph{Official setting.} As in Cell~1 (\S\ref{p2c1:sec:count}): $U(G)$ is the least number
of uphill paths over all labellings of $G$. Values are handed in together with an explicit
labelling, listed as the $2^d$ vertices in increasing label order.

\begin{statement}
Determine $U(Q_5)$, and give a labelling attaining it.
\end{statement}

\paragraph{Status.} \textbf{Solved}. The upper bound is by hand. The lower bound is
\textbf{computer-assisted}: one SAT computation, below one second, repeated with two solvers.

\paragraph{Answer.} $U(Q_5)=88$, and $\nabla(Q_5)=14$ (Theorem~\ref{p2c2:thm:main}).
The labelling is in \S\ref{p2c2:sec:lab}.

\subsubsection{The decycling number of $Q_5$}

The counting bound of Lemma~\ref{p2c1:lem:count} gives only $4\nabla(Q_5)\ge3\cdot16+1$, that is,
$\nabla(Q_5)\ge13$, and so $U(Q_5)\ge84$. This bound is not attained.

\begin{lemma}[computer-assisted]\label{p2c2:lem:L5}
$Q_5$ has no decycling set with $13$ or fewer vertices. Hence $\nabla(Q_5)\ge14$.
\end{lemma}

\begin{proof}[Method, and why the computation proves the lemma]
This is claim (L5) of \texttt{code/uphill\_lower.py}. The variables are $x_v$ ($v\in V(Q_5)$),
with $x_v$ meaning $v\in R$. The formula $F$ contains:
\begin{enumerate}[label=(\alph*)]
\item a CNF encoding (sequential counter) of $\sum_vx_v\le13$;
\item for every $4$-cycle $K$ of $Q_5$, the clause $\bigvee_{v\in K}x_v$;
\item the clauses $\bigvee_{v\in K}x_v$ for cycles $K$ found during the loop below.
\end{enumerate}
Each clause of type (b) or (c) says that $R$ meets the cycle $K$, and every decycling set does
this. So every decycling set of size at most $13$ satisfies $F$ at every stage.

\emph{Loop.} The solver returns a model $R$, or reports UNSAT. If it returns $R$, we compute
$Q_5-R$ with \texttt{networkx}. If this graph has a cycle, we add the clauses of a cycle basis of
every cyclic component (type (c)) and ask again. The new clauses are violated by the current
model, so no model is returned twice. Hence the loop ends, after finitely many rounds. It ends in
one of two ways: a model $R$ for which $Q_5-R$ is a forest (a decycling set of size $\le13$), or
UNSAT.

\emph{Result.} UNSAT, after $348$ rounds in $0.4$\,s with CaDiCaL~1.5.3, and after $357$ rounds in
$0.3$\,s with Glucose~4, both through \texttt{python-sat}. Since every decycling set of size
$\le13$ would be a model, none exists. The only facts used are the correctness of the
cardinality encoding and of the SAT solvers. Running two independent solvers guards against a
solver bug.
\end{proof}

\begin{theorem}\label{p2c2:thm:main}
$\nabla(Q_5)=14$ and $U(Q_5)=88$.
\end{theorem}

\begin{proof}
Lower bound: Lemma~\ref{p2c2:lem:L5} and Theorem~\ref{p2c1:thm:lower} give
$U(Q_5)\ge32+4\cdot14=88$. Upper bound: $C=\{00000,11110\}$ lies in $E_5$ and has distance $4$.
By Proposition~\ref{p2c1:prop:code}, $E_5\setminus C$ is a decycling set of size $16-2=14$, and
the labelling $L_C$ has $32+4\cdot14=88$ uphill paths.
\end{proof}

\begin{remark}
A set $C\subseteq E_5$ with pairwise distances $\ge4$ has at most two elements. Suppose it had three.
Translating by one of them (an even vector, which preserves $E_5$ and all distances), we may assume
$C=\{0,a,b\}$. Then $|a|=|b|=4$, so $a$ and $b$ are $4$-subsets of a $5$-set and meet in at least
$3$ coordinates. This gives $|a\oplus b|\le2$, a contradiction. So Proposition~\ref{p2c1:prop:code}
cannot give less than $88$, and Lemma~\ref{p2c2:lem:L5} shows that no labelling can.
\end{remark}

\subsubsection{The labelling}\label{p2c2:sec:lab}

$L_C$ with $C=\{00000,11110\}$: codewords, then the $16$ odd vertices, then the remaining $14$
even vertices, each block in increasing binary order. It has $88$ uphill paths (checked with
\texttt{code/uphill\_eval.py}).
\begin{flushleft}\small\ttfamily
\makebox[2.4em][r]{\textrm{1:}}\ 00000\ 11110\ 00001\ 00010\ 00100\ 00111\ 01000\ 01011\\
\makebox[2.4em][r]{\textrm{9:}}\ 01101\ 01110\ 10000\ 10011\ 10101\ 10110\ 11001\ 11010\\
\makebox[2.4em][r]{\textrm{17:}}\ 11100\ 11111\ 00011\ 00101\ 00110\ 01001\ 01010\ 01100\\
\makebox[2.4em][r]{\textrm{25:}}\ 01111\ 10001\ 10010\ 10100\ 10111\ 11000\ 11011\ 11101\\
\end{flushleft}
